

\hspace{-5em}\begin{minipage}{1.1\textwidth}
\begin{flushright}
{\fontsize{16}{22}\selectfont
\textbf{%
\dunderline[-8pt]{3pt}{PROLOGUE}}
}
\end{flushright}
\end{minipage}


\vspace{6em}


%\hfill

\hspace{-12.5em}\begin{minipage}{1.3\textwidth}
\begin{flushright}
{\fontsize{19}{22}\selectfont
\textsf{\textit{What Has Changed in Two Decades, and What\\\vspace{5pt}
Remains the Same?}}
}
\end{flushright}
\end{minipage}

\vspace{9em}


\prSect{INTRODUCTION}

\pfl{}

\noindent{}We begin our updated critique of the American family court system by
acknowledging that, notwithstanding the fact that mothers are often
placed at risk of incarceration, fines, and losing all contact with their
children in retaliation for requesting the family courts to shield their
children from abuse committed by the other parent, there have still been some successes in this uphill battle to reform the family courts. Such
successes give us reason for hope. By elucidating the problems prevalently found in the family court system, we see evidence of steady efforts
to effect change. Because for those who have become aware of the troubling patterns of family court dysfunction, it becomes increasingly hard for them
to accept the dismal status quo. \pfl{}

Most notably, a cadre of reformers have set the wheels in motion for
course-correcting a system that has careened out of control in those\mdash{}all
too frequent\mdash{}instances in which family court decisions are contrary 
to the safety and well-being of abused
children. Taking a balanced approach, we acknowledge the fact that
there are many good judges who zealously protect children from harm.
However, there are many who do not. And too many children in custody
matters are thus falling through the cracks in the system. \pfl{}

Among those who have had a significant effect on improving the
family court and child welfare system is Dr. Jeffrey L. Edleson, whose
noteworthy contributions will be discussed at length in the addendum to
%\notePage{xxxiii}
the third section of this book. Specifically, Dr. Edleson, in collaboration
with his colleague, Susan Schechter, authored a comprehensive guide
for court professionals and child advocates in protecting the welfare and
safety of children. Their seminal book, \textit{Effective Intervention in Domestic
Violence and Child Maltreatment: Guidelines for Policy and Practice}\fnm{1}\mdash{}
better known as the \qq{Greenbook} (after the color of its cover)\mdash{}whose
focus is \qq{to present leaders of communities and institutions with a
context-setting tool to develop public policy aimed at keeping families
safe and stable,}\hspace{1pt}\fnm{2} has spawned an impressive body of literature
based on its comprehensive and cogent recommendations for guiding
professionals in working with battered women and abused children. \pfl{}

Although it was first written in 1999, a concrete example of the effects of
the Greenbook in improving child safety appeared twenty years later
in a 2019 journal article. In \qq{Advocacy Perspectives on the Greenbook,}
Lucy Salcido Carter explained that \qq{[w]hen domestic violence is an issue
in the case, it is important not to have [issued] an automatic finding
of failure to protect against the non-offending parent [the mother], but
instead to recognize that abuse is happening to the adult victim and that
the perpetrator is exposing the child to the domestic violence.}\hspace{1pt}\fnm{3} \pfl{}

Similar to the Greenbook's introduction of helpful policies and practices in protecting children, on the clinical front some noteworthy
advances have been made that deserve mention. Most notable, eminent neuropsychologist and family counseling expert Dr. Robert (Bob)
Geffner spearheads a unique continuing education training program
for clinicians that introduces to the participants advanced research and clinical methods
for assessing and treating intrafamilial abuse and trauma. Founder of
the Institute on Violence, Abuse and Trauma (IVAT),
Geffner is responsible for providing thousands of therapists with training in abuse and trauma intervention since IVAT's creation in 1984. With a focus on educating the mental health profession about the traumatic effects on children judicially ordered into the custody of a sexually abusive parent, 
Geffner has paved the way for children and adult survivors to receive the trauma intervention therapy they need to heal from
exposure to ongoing assault. His work has resonated throughout the
mental health profession and has established a new standard. \pfl{}

%\notePage{xxxiv}
All in all, dedicated researchers and clinicians have taught us that it requires patience and forbearance to rehabilitate a system that has become
tainted with prejudices, myths, and bad practices. We hope this
painstakingly candid critique of the family court system will motivate,
inspire, and galvanize change in how judges approach a custody decision 
involving a protective parent, where the other parent is credibly accused
of child sexual abuse, so that tendentious mental health theories 
interlaced with gender, racial, ethnic, and social class biases no longer inform
destructive family court rulings, as they so often have. \pfl{}

We must not forget that without an enlarging contingent of family
court reformers, our successes are likely to remain limited and 
fragmentary.  There have been small but measurable successes\mdash{}more 
judges, and others in the family court system, have been made aware of the pitfalls
of system biases that can endanger children's safety and 
welfare\mdash{}but
we continue to observe serious wrongs. In case after case, the evidence
reveals a system that all too often turns its might against protective
mothers and consigns abused children to a tragic future of helpless suffering. 
No wonder so many mothers caught up in family court litigation
describe their experience as \qq{trial by ambush.} In case after case, protective 
parents face an unremitting onslaught of legal maneuvers, all
too often involving aggressive collusion between the allegedly abusive
parent, caseworkers from Child Protective Services (CPS), and even the
child’s court-appointed attorney. And abuses like these are particularly
harmful to Black, Indigenous, and People of Color (BIPOC) litigants, who
have to shoulder the burden of race on top of the systemic biases and
shortcomings of the family court system. \pfl{}

For these reasons, we offer a second edition of \textit{From Madness to
Mutiny} in the hope that it will stimulate purposeful and productive 
discourse\mdash{}not only in the classrooms that will produce the next generation of lawyers,
judges, and social workers, but in the hospitals, medical offices, 
and child evaluation centers where evidence of domestic violence and child abuse 
is a recurring reality. Sociologists, psychologists, feminists, policy makers,
%\notePage{xxxv}
and others must be enlightened to the scourge of gender bias that has
entrapped far too many mothers in the family courts. Those mothers suffer the anguish of seeing their children harmed; many are subsequently
expunged from their children's lives or restricted to minimal contact
with them. We believe that this malady
must be attacked at its roots. This was the purpose of the first edition of
our book, and likewise the warrant for the second.\fnm{4} \pfl{}


\prPart{PART{\prS}I---A{\prS}FAMILY{\prS}COURT{\prS}CUSTODY%
{\prS}CRISIS{\prS}OF{\prS}EPIC\\PROPORTION}

%\prSect{PART I---A FAMILY COURT CUSTODY CRISIS OF EPIC PROPORTION}


%\subsect{15}{United{\prSb}Nations{\prSb}Addresses{\prSb}Family{\prSa}Courts'{\prSa}%
%Pervasive{\prSa}Gender\\\vspace{-1pt}Bias{\prSs}and{\prSs}
%the{\prSs}\qq{Double}{\prSs}Victimization{\prSs}of{\prSs}Women{\prSs}%
%and{\prSs}Children\\as{\prSa}a{\prSa}%
%Global{\prSa}\makebox{Human{\prSa}Rights{\prSa}Crisis}\par}

\subsect{15}{United{\prSb}Nations{\prSb}Addresses{\prSb}Family{\prSb}Courts'{\prSb}%
Pervasive{\prSb}Gender\\\vspace{-1pt}Bias{\prSs}and{\prSs}
the{\prSs}\qq{Double}{\prSs}Victimization{\prSs}of{\prSs}Women{\prSs}%
and{\prSs}Children\\as{\prSa}a{\prSa}%
Global{\prSa}\makebox{Human{\prSa}Rights{\prSa}Crisis}}


\noindent{}In the fall of 2022, the United Nations Human Rights Office of the
High Commissioner issued a \qq{Call for Inputs\mdash{}Custody Cases, Violence
against Women and Violence against Children,} which was a far-ranging 
request to academic institutions, international organizations, 
and governmental agencies for data pertaining to the gravity of abuses perpetrated
against women and children caught in the family courts around the
world. \pfl{}


Pointing to family courts' tendency to dismiss evidence of domestic
violence and child abuse in custody cases, the \qq{Call for Inputs} asserted
that many women who are victims of violence and whose children have
been abused by the father will \qq{face double victimization as they are
punished for alleging abuse,} and by \qq{losing custody} of their children
\qq{or at times being imprisoned.} Identifying the crux of these
disastrous custody decisions, the UN Office of the High Commissioner
pointed to family courts' widespread reliance upon concepts such as
\qq{parental alienation} (in which a parent, usually the mother, 
is unjustly blamed for any hostility shown by an allegedly abused child toward its
father). The High Commissioner's Office
pointed out that though \qq{these concepts lack a universal clinical or scientific definition,} they are widely accepted in the family courts as 
justification for stripping a mother of custody of her child. \pfl{}

%\notePage{xxxvi}
Driving this major UN Human Rights project was a need \qq{to examine the ways in which family courts in different world regions refer
to parental alienation, or similar concepts, in custody cases and how
this may lead to double victimization of victims of domestic violence or
abuse} and \qq{to document the many ways in which family courts ignore
the history and existence of domestic \ellThree{1pt}{-1pt} violence and [child] abuse in
the context of custody cases} while studying \qq{their grave consequences
[to] mothers and their children.} \pfl{}

Furthermore, the Human Rights Office of the High Commissioner
pointed out that one of the goals of their international study is to learn
more about \qq{any intersecting elements such as age, sex, gender, race,
ethnicity, legal residence, religious or political belief} that result in poor
case outcome for women and children. They stated the overall purpose
of their study is to \qq{draw attention to the scale and manifestation} of the
double victimization of women and children when courts across \qq{many
countries, spanning all regions of the world} discredit mothers' pleas
to protect their children from abuse and to \qq{offer recommendations for
States and other stakeholders to address the situation.}\hspace{1pt}\fnm{5} \pfl{}

The UN study offered hope to women's groups across the globe. In
the United States, The Women's Coalition\mdash{}a well-organized group of
mothers and advocates devoted to collectively raising awareness about
the custody crisis\mdash{}was optimistic about the project. The group's founder,
San Diego resident Cindy Dumas, a PhD candidate in psychology whose
career was derailed when she became a victim of the family court and
was forced to go on the run with her children, posted the following on
her website:

\prQuote{This is the first time a mechanism of the UN HRC has prioritized this issue and initiated an inquiry. It is a leap forward in our
efforts to get it officially recognized as a human rights violation.
The Women's Coalition has previously submitted three petitions,
which have likely contributed to them initiating this inquiry.\fnm{6}} \pfl{}

Shortly after the UN's announcement of its \qq{Call for Inputs,} The Women's Coalition 
distributed a questionnaire to protective mothers.
%\notePage{xxxvii}
In the first two weeks, Dumas received nearly four hundred
completed questionnaires from mothers who had lost their children to an
unsympathetic family court system, many of whom having lost not 
only contact but visitation rights with their children, 
from whom (in some cases) they were now permanently estranged.
Her questionnaire covered all
the component parts of the family court system that contribute to the
cruelest rulings against mothers. But even more so, it captured the harm
inflicted on a child who is forced to live with an abusive father\mdash{}
harm that often includes induced hostility against the mother so strong that
reconciliation becomes impossible. \pfl{}

By early summer 2023, the UN's Human Rights Office of the High
Commissioner would produce a conclusive report on its findings. The
press release accompanying its report had an alarming heading:
\qq{Urgent reforms needed to protect woman and children from violence
in custody battles: UN expert.} The report used harsh language 
in its description of
the family courts: \qq{Deeply embedded gender bias that pervades family court
systems across the globe is placing women and children in situations of
immense suffering and violence, a UN expert said today.} \pfl{}

UN Special Rapporteur Reem Alsalem, who conducted the study,
went on to state that \qq{the tendency of family courts to dismiss the history of
domestic violence and abuse in custody cases, especially where mothers
and/or children have brought forward credible allegations of domestic
abuse, including coercive control, physical or sexual abuse is unacceptable.} Perhaps most telling was Alsalem's warning about the
effects of a failed family court system on BIPOC mothers: \qq{The report
underscores that minority women face additional barriers when being charged of 
[\textit{sic}] using \sq{parental alienation} in part due to increased
barriers in accessing justice as well as negative stereotypes.}\hspace{1pt}\fnm{7} \pfl{}

\vspace{3em}
\subsect{16}{Inside the Torture Chambers of the Family Courts\raisebox{2pt}{:} Forced
Weaning of Young Infants to Accommodate Overnight Visits
to the Abusive Father}

\noindent{}In recent years, family courts have removed very young children,
including nursing infants, from mothers with much greater frequency
%\notePage{xxxviii}
than in the past. Even twenty years ago, when researching the first edition of
this book, we found cases in which infants and toddlers were removed
from protective mothers. However, today’s family courts are still more
likely to take such drastic steps. \pfl{}

Such orders have consequences that can border on the grotesque. A
nursing mother, though forbidden to keep her child overnight, may be
forced to deliver bottles of her expressed breast milk to an allegedly 
abusive father so that \textit{he} can hold the child at his mercy for as long as the
family court deems proper. \pfl{}

Even when the court’s order requires \qq{only} that the mother deliver
the child to the allegedly abusive father for three or four overnight visits
each week, the consequences can be harrowing. Some nursing infants do
not adjust readily to expressed milk (or formula), in which case they may
go hungry, which the mother will have to deal with after she retrieves
the child. A mother caught in this bind may try to accustom the baby
to bottled milk or formula instead of breast milk solely to accommodate
the court’s directives. But while this prevents the child from being 
dehydrated while in the father's custody, it is probably not in the child's best
interests. Notably, both the World Health Organization (WHO) and the
American Academy of Pediatrics recommend breastfeeding of an infant
for at least one year, and ideally for two to two and a half years. 
The reason for this is that breastfeeding inures to the benefit of the baby both emotionally and physically
by strengthening its immune system to fight off disease, improving its
growth and development, and fostering a healthy maternal-child bond. \pfl{}

In the winter of 2023, a Virginia mother made headline news when
the judge forced her to wean her four-month-old infant so that the infant's father could have extensive visitation with the child, even though
he had abandoned the mother during her pregnancy.\fnm{8} \textit{The Washington
Post}, which broke the story, set off a firestorm when the fathers' 
rights activists falsely claimed the mother was using breastfeeding as a \qq{weapon
against visitation.} One of the authors of this book (Amy Neustein)
spoke to the mother at length and learned of her Sisyphean struggle
against a court system that that had condemned her 
for being \qq{uncooperative.} he record of her case did not support the court's criticism;
in fact, there was evidence that the mother had gone out of her way to
encourage the father’s participation in the life of their child. Moreover,
she had offered to have him come to her home for visits with the child,
so that she could nurse the baby as needed during the visits. The court
required her, instead, to express her milk so that the father could spend
his visits elsewhere. \pfl{}


Concerned about the mistreatment of this lactating mother,
Neustein\mdash{}together with author Michelle Etlin\mdash{}wrote a letter to 
\textit{The Washington Post}. The paper published their letter under the heading
\qq{Infants are being penalized in the name of \sq{men's rights.}} Their letter
begins with a sad but true commentary: \qq{This story is all too familiar to
those who study the family court system. For the past four decades, we
have researched how U.S. domestic relations courts\ellThree{1.5pt}{2pt}have cruelly and
senselessly torn infants and children from their mothers.}\hspace{1pt}\fnm{9} \pfl{}


\vspace{3em}
\subsect{15.5}{Another{\prSb}Instance{\prSb}of{\prSb}Violence{\prSb}%
against{\prSb}Women:{\prSb}Jailing{\prSb}of{\prSb}\\Mothers%
{\hspace{7pt}}on{\prSc}Untenable{\prSc}Grounds}

\vspace{-5pt}
\noindent{}Another problem that seems to be on the rise in family courts is the 
arbitrary jailing of mothers. One protective mother in Maryland, who is
diabetic, told us that she was jailed for a day by a family court judge
who was angry because, due to a diabetic episode, she was having
trouble attending to his questions. Another mother\mdash{}this one in New
York's Westchester County\mdash{}was jailed for arriving late for a hearing,
even though she explained to the judge that she had been delayed in
bumper-to-bumper traffic. \pfl{}

\enlargethispage{25pt}Mothers, even if not jailed, have in recent years been threatened with
incarceration for talking to their own supportive family members about
their family court proceedings, even if the family members will never
need to serve as witnesses either for the mother or as \qq{hostile} witnesses,
called by the other side. Forbidding a desperate mother from seeking
moral support from her own family members causes her to feel isolated,
demoralized, and browbeaten. Some mothers become suicidal, while
others have actually succeeded in taking their own lives. \pfl{}\newpage{}

\subsect{15}{Ending Institutional Violence against Mothers and Children}


\noindent{}The sum total of such oppressive measures\mdash{}punitive deprivation of
child custody, judicial interference with breastfeeding, and unjustified
incarceration\prsa{}of\prsa{}family\prsa{}court\prsa{}litigants{\prsc}\mdash{}{\prsb}%
represents\prsa{}a\prsa{}genuine\prsa{}human\\
rights crisis, and not to mention\mdash{}in many cases\mdash{}violations of law and
of constitutional protections. This senseless infliction of institutional 
violence on women and children demands a thorough investigation by the
United States Attorney General. And such an undertaking may be less
unlikely than it sounds.  \pfl{}

Legal scholar and first champion of the protective mother, Jeremiah
B. McKenna, to whom we have dedicated this book, concluded after
research into the relevant federal statutes, that a US Department of 
Justice investigation into family courts and their collaborating auxiliaries
(Child Protective Services, law guardians,\fnm{10} foster care agencies, 
visitation centers, etc.) was not only feasible but entirely justified by the facts.
Based on McKenna’s research Amy Neustein directed a thirty-one-page
memorandum to then-United States Attorney General Merrick Garland
(under the Biden Administration) describing family court abuses in
a few particularly dysfunctional counties.\fnm{11} A federal prosecution can
accomplish more than the application of criminal penalties to a few individuals. 
It can signal to all family court officials and auxiliaries that the
corrupt practices that mar so many child-related cases\mdash{}rejection of 
relevant evidence, removal of children without due process of law, denial
to mothers (and especially mothers who are women of color) of basic
civil rights\mdash{}will no longer be tolerated. \pfl{}

In fact, the paradigm on which the existing system was designed
may prove unworkable altogether. It may, given its deep-seated venality,
become defunct, to be replaced perhaps by a panel of concerned 
professionals and lay members who are not driven by any agenda, as their only
goal would be to protect the mother-child bond and shield the child
from predatory harm. This would be more akin to a \qq{King Solomon's
Court}\mdash{}where wisdom and compassion would determine the fate of a
child rather than politicking, cronyism, and political favors. \pfl{}


\prPart{PART II\mdash{}PROTECTIVE MOTHER MOVEMENT: THEN AND NOW}


\subsect{16}{Protective Mother Movement Makes Great Strides in Recent Years}


\noindent{}In the nearly two decades since the publication of the first edition of this
book\mdash{}which showed how courts too often punish mothers with the
loss of custody of their children when they make good faith allegations
of sex abuse against the child's father\mdash{}we have witnessed we have witnessed the burgeoning
of a social movement to reform the family courts.  This is due in part to the elaborate networking offered by social media and, more recently, by Zoom
%\notePage{xlii}
meetings that bring together protective mothers across the country for
an informal but important virtual meetings. There are now scores of trusted websites (and blogs) posting the latest relevant research studies, news articles, and
video and YouTube recordings, all available to mothers who have lost
their children to the family courts. \pfl{}

Hence, no mother is fighting alone today\mdash{}isolated in her despair\mdash{}
for she can now readily avail herself of the resources available from a
multitude of websites, where some of which provide referrals to family 
law attorneys, potential expert witnesses, or give advice on legal
strategies for winning back custody of children. Some of these websites
propose specific legislative reforms to improve the functioning of the
family courts while others propose a jettisoning of the existing court
system altogether, to be replaced by a more humane structure for de-
termining children's best interests. Regardless of their approach, the
mantra of all these sites\mdash{}National Safe Parents Organization,\fnm{12} Mothers ReVolution,\fnm{13} Custody Peace,\fnm{14} The Women's Coalition,\fnm{15} 
Safe Kids International,\fnm{16} One Mom's Battle,\fnm{17} 
The Leadership Council,\fnm{18} and California 
Protective Parents Association,\fnm{19} just to name a few\mdash{}is to raise
public awareness about the catastrophically high numbers of mothers
losing custody to abusive ex-partners or spouses. \pfl{}


The term \qq{protective mother}\mdash{}which has become a part of the vernacular of family law practitioners and women's groups across the
country\mdash{}was coined around 1986 by mental health professionals who
had formed the American Professional Society on the Abuse of Children
(APSAC). It refers to a mother who mounts a fervent struggle to protect
her child from all forms of abuse, particularly sexual abuse committed by 
the father, within a hostile family court system intent on beating her
down every step of the way. And we would add that another
part of the protective mother's struggle is her engagement in organized
protests against the broader effects of family court inequities. Some of
those activities have included participation in peaceful demonstrations
in front of state capitols around the country as well as at the nation's
capitol. Letter-writing/e-mail campaigns, petitions, and social media
posting/likes have served as another way of voicing concerns about the
crisis in the family courts. A five-minute Google search of the protective
parent/family court issue will provide an excellent gauge of the rhythm
and beat of this growing social movement. \pfl{}

\vspace{3em}
\subsect{15}{Plight of Protective Mothers Heeded in Academic Settings}

\noindent{}The protective mother issue now resonates in the halls of the universities, whereas two decades ago when the first edition of this book was
being prepared for publication there had only been sporadic interest in this 
topic at the institutions of higher learning. For example, one of the authors 
(Amy Neustein) was invited
in the fall of 1990 to lecture at the Women's Law Forum at Rutgers State
University School of Law\fnm{20} only because the professor organizing this
effort, H. Joan Pennington, was a remarkable activist and visionary, having founded the National Center for Protective Parents, Inc., as a legal
resource for mothers caught in hostile family court litigation. But such
efforts to bring this subject into undergraduate and graduate classrooms
at that time were certainly the exception to the rule. \pfl{}

Today, the Battered Mothers Custody Conference, co-chaired by Professor Mo Therese Hannah, is held annually at Siena College in Albany,
New York, where Professor Hannah serves as a full professor of psychology.
The attendance at her conference has skyrocketed since
its inception in 2004; it now serves as the premier nexus for the 
protective mother movement. More recently, the California Protective Parent
Association has partnered with the University of California Irvine (UCI)
Initiative to End Family Violence; and in doing so, this grass roots
parents' group has joined a major academic institution to collaborate
on an annual, full-day, accredited continuing legal education 
symposium (although at present, held virtually), which started in fall 2021. The
conference, titled \qq{Forward Together: Multidisciplinary Perspectives on
Protecting Children from Abuse,}\hspace{1pt}\fnm{21} had a panel of speakers consisting
of several activists who addressed the colossal mistreatment of mothers
in the family court system. All of this points to a sea change in attitude. \pfl{}

\vspace{-1.4em}
That is, whereas mothers' disenfranchisement in the family courts\mdash{}
which had at one time been solely the province of activists\mdash{}had now
entered mainstream institutions for discussion and debate. And as
%\notePage{xliv}
previously mentioned, the discourse on this topic was not a flash in the
pan; instead, it has become a regularly occurring annual event. This is
a marked change from 2005 when the first edition came out. Back then,
Amy Neustein had been invited to serve as visiting faculty at the
Tristate Family Law Conference held by the National Judicial College
in Bretton Woods, New Hampshire.\fnm{22} The attendees at the conference,
and the conference organizer herself, looked upon this subject as novel,
but no less worthwhile to explore. This represented a healthy discourse
on family courts' betrayal of women and children inside the protected
corridors of a judicial education institution, although back then such a
forum was more or less an isolated event. \pfl{}

As the years moved onward, and as academic institutions began
to show a continuing and steadfast interest in the atrocities of family court, George Washington University Law School would become
another very important venue for broadcasting these profound injustices toward mothers who tried to protect their children from abuse.
George Washington University Law School Clinical Law Professor
Joan S. Meier became the Founding Director of the National Family
Violence Law Center (NFVLC) situated at George Washington University Law School. 
The center's function is described as follows: \qq{using empirical research,
policy development, and collaborative advocacy the NFVLC
educates professionals, the media, and the public about the problems
in court responses to adult and child abuse, and advances critical legal
reforms and systemic change.}\hspace{1pt}\fnm{23} \pfl{}

\enlargethispage{1em}%
In fall 2021, \textit{Family and Intimate Partner Violence Quarterly} published
a study, titled \qq{Protective Mothers, Endangered Children,} spearheaded 
by Professor Geraldine Butts Stahly of California State University at
San Bernardino, together with two highly respected activist organizations,
the California Protective Parents Association and Child Abuse Solutions, Inc. 
The data collected for this study came from a hard-copy
questionnaire, consisting of 101 questions, which had been distributed
at national and international conferences concerning domestic violence
and child abuse. The purpose of this study was to learn about the \qq{experience of mothers attempting to protect their children from abuse.}\hspace{1pt}\fnm{24}
The authors of this study were so alarmed by the results that they began
%\notePage{xlv}
to circulate its findings (399 mothers from across the country had filled
out the questionnaire) at multiple workshops and poster sessions \textit{ahead
of} its publication as they felt an urgency to alert professionals to the
serious failures of the family courts. \pfl{}\newpage{}

Specifically, the study authors observed evidence that, again and
again, mothers lost custody of children to fathers who appeared to
place them at risk. In California, where the rate of such system failure
appeared to be especially high, mothers lost custody of children to 
identified abusers in 85.4 percent of the cases described, as opposed to 74.6
percent in other states. The study authors asked for an overhaul of every
facet of the system\mdash{}custody evaluators, mediators, judges, child 
welfare services, attorneys for children, state law enforcement\mdash{}including
setting up \qq{a citizen review \sq{second look} 
commission \ellThree{-1pt}{-1pt} at the state
and federal level to review cases in which children are currently at risk,
with the power to intervene [in order] to protect [children].} The study
concluded with an exhortation for \qq{[i]mmediate and effective changes
[that] must be made to law, policy and procedure to provide safety for
children of divorce and separation.}\hspace{1pt}\fnm{25} \pfl{}

\vspace{3em}
\subsect{15}{Protective Mothers Finally Find a Sympathetic Ear in the Mass Media}

\noindent{}On the media front, likewise, we have witnessed significant progress in
addressing the protective mother issue since we first wrote about it in
2005. Those changes in the news media's posture deserve mention. To
wit, in the early days of the protective mother movement, the media
coverage was sporadic at best. But even more troubling was the fact that
it was so unpredictable at times, and often biased against the mother. \pfl{}

As a result, many mothers were \qq{duped} when reporters assured them
the hard news coverage or feature story would be in their favor, only
to learn too late of the influence of fathers' rights groups on the press at
that time. Indeed, quite a few protective mothers, having struggled to
convey their stories to reporters, found themselves portrayed in a 
terribly unflattering light when the report finally appeared. 
By contrast, the
fathers they suspected of abusive conduct, and even more so their 
advocates\mdash{}including spokesmen for fathers' rights organizations\mdash{}usually
appeared as respectable and credible, regardless of the facts of the case. \pfl{}

%\notePage{xlvi}

More often than not, the copious documentation supplied by the
mothers to prove their innocence had no effect on the outcome of the
news media story. However, since the publication of the first edition,
we have seen more thoughtful and accurate reporting of these issues; it
is less common to find myths and prejudices about protective mothers
reported as fact. \pfl{}


For example, in 2016, investigative journalist Laurie Udesky, working for 100Reporters, a nonprofit news organization giving a platform
for freelance journalists writing about corruption and social injustices,
published a story on the plight of mothers in the family courts. 
Her article was titled \qq{Custody in crisis: How family courts
nationwide put children in danger.}\hspace{1pt}\fnm{26} A year later she would be given
the Excellence in Journalism Awards, from the Society of Professional
Journalists, Northern California Chapter, for her work in exposing the
ill-fated decisions of family courts that place children in harm's way. \pfl{}

In 2019, \textit{The Washington Post} would write a feature story on the abject plight of mothers punished with the loss of custody when they
try to protect their children from abuse perpetrated by the other parent. 
The article, aptly titled \qq{\hspace{2pt}\sq{A gendered trap,}} quoted Meier,
who found a pattern of judges repeatedly placing children with the
abusive parent: \qq{Meier called the pattern a \sq{complete betrayal of the
mission} of the courts, which are supposed to prioritize the welfare of
children.}\hspace{1pt}\fnm{27} \pfl{}


What prompted that prominent news coverage in \textit{The Washington Post}
was Meier's large-scale study that found that \qq{mothers who report
abuse\mdash{}particularly child abuse\mdash{}are losing child custody at staggering
%\notePage{xlvii}
rates.} In fact, the study showed that in over 73 percent of the time when
mothers alleged abuse, and the other parent cross-claimed \qq{parental
alienation} (the notion that when a child refuses to have a relationship
with a parent it is due to purported \qq{manipulation} by the other parent), the family court judges will switch custody to the alleged abuser.\fnm{28}
The Meier study, as a matter of fact, was able to hold up admirably to
specious attempts made by two critics\mdash{}with a clear agenda\mdash{}to
negate these significant study findings. \pfl{}

Namely, in 2022, the \textit{Journal of Family Trauma, Child Custody, and
Child Development} (formerly the \textit{Journal of Child Custody}) published a
strong and persuasive commentary by Meier and colleagues that pointed
out that opponents of her study (although only two\mdash{}certainly not
a huge number) had based their criticism on \qq{incorrect speculations
and misstatements about [their] methodology} and that any possible
\qq{methodological flaws} that may have \qq{legitimately cast doubt on the
validity of the FCO [Family Court Outcome] study's findings have not,
to date, been identified.} In responding to her critics, 
Joan S. Meier and her coauthors put 
this matter finally to rest by showing the \qq{lack of substance}
to their critics' findings \qq{while substantiating the credibility} of their assiduous research of family court custody (case) outcome when parental
alienation is raised to discredit the mother.\fnm{29} \pfl{}

In her successful rebuttal of her critics' fallacious arguments, Meier
preserved the important contributions of her study and, no less, the
eye-opening lead story in \textit{The Washington Post}. This was crucial, as any
progress made in getting exposure for the protective mother issue in a
major outlet, such as \textit{The Washington Post}, might have been reversed if
the critics of her study have gone unanswered. What we see here is that
the deeply rooted gender biases casting aspersions on the credibility of
mothers cannot be left unchallenged in either the professional circles\fnm{30}
or in the mass media. As for the latter, because it is axiomatic that public
opinion drives policy makers and legislators, it is imperative to have 
media coverage that is supportive of protective mothers instead of working
against them. \pfl{}


While the 2019 \textit{Washington Post} article still remains the most highly
cited news story today about the undeserved backlash against mothers
%\notePage{xlviii}
who in good faith report child sexual abuse to the courts, a handful of
mainstream media venues have, in subsequent years, reported on this
topic as well. For example, in May 2021, NBC aired a two-part series
on children who were murdered by their fathers during unsupervised
contact with their children. Pointing to the lack of judicial oversight or
accountability, child advocates who were interviewed for this NBC series 
argued for an overhaul of the court system, including the passage
of laws that would compel courts to consider abuse evidence\mdash{}both
present and past\mdash{}when making custody/visitation decisions so as to
ensure children are protected from violence and abuse and, in the most
extreme case, death.\fnm{31} \pfl{}


Two years later, in 2023, New Hampshire investigative journalist
and author Olivia Gentile wrote an outstanding feature-length story
for \textit{Insider} on how the discredited \qq{Parental Alienation} theories have
contaminated the American family courts, causing abused children to
be confined to the custody of their abuser with little or no access to
their mothers, who strenuously try to protect them. 
Reporting on a San Diego
case, as an example of this phenomenon, Gentile showed how a mother
lost custody of her children after her son's therapist reported to the
court how the twelve-year-old boy was coerced by the father to 
perform lewd acts that the child described as \qq{disgusting.} In her reporting
on this story, Gentile showed that several years after the mother lost
custody of her three children to their abusive father, the mother 
accidentally found in a cloud storage account she had once shared with
her ex-husband numerous sexually explicit images (both video and still
shots) the father had taken of the children after he obtained custody of
them.\fnm{32} \pfl{}

In our interviews with women who had conversations with reporters,
we found that journalists like Gentile were deeply appreciated by the
protective mothers. What we heard was that she gave generously of
her time to distressed mothers, often answering their e-mails\mdash{}often at
late hours of the night, and on weekends and holidays, too. Notably such
reporters were not feminist writers with a long history of championing
women's causes; instead, they were sensitive and caring 
journalists who witnessed an injustice meted out 
to protective mothers and applied their
journalistic skills to try to bring these family court 
travesties to light. \pfl{}


\vspace{2em}

\subsect{15}{Feminist Writers and Researchers Come Together to Champion
the Plight of Protective Mothers to a Broader Audience, Outside
of the Traditional Feminist News Outlets}

\noindent{}This past decade has proved historic for broadcasting the protective mother issue to a wider readership. Whereas in the early days,
namely the 1980s, 1990s, and early 2000s, it was mainly feminist newsletters and magazines that had profiled the wretched fate of mothers
peregrinating through the family court and child welfare system, today
one can find discussion of these critical issues in venues outside the feminist media. And the researchers who have labored to study this problem
have much to do with the progress in that regard. \pfl{}

For example, \textit{Providence Journal} columnist Anne Grant, who is the
former executive director of the Women's Center of Rhode Island, has
paid much attention to the gender biases pervasive in the family courts
in her writings on social inequities and unbalanced power across various sectors of society. The \textit{Providence Journal} is certainly not a feminist
publication as it is a general-purpose newspaper widely read in the New
England area. In her Amazon review of \textit{Beyond the
Hostage Child: Toward Empowering Protective Parents}\fnm{33}
by Leora N. Rosen, PhD (former senior social science 
analyst at the National Institute of Justice), Grant wrote:

\vspace{-1em}
\prQuote{This book holds validation for those who have been traumatized
when courts removed terrified children from protective parents
and gave them to the sole custody of abusers. Dr. Rosen shines
a light we need to go forward. \quasiPar{}
She asserts that alleged crimes of domestic violence and child
sexual abuse within the family should never be sent to civil
courts that are designed for compromise. She briefly describes five
%\notePage{l}
proposed models for change and offers more detail on a sixth,
composite model, Child At Risk Classification Office (CARCO)
that focuses on a public health assessment of the child's risk of being exposed to violence or abuse. She uses the acronym TRIAL to
represent key elements of CARCO: Training, Reporting, Investigation, Adjudication, and Long-term planning\\\mdash{}that are woefully
absent from the present practice of adversarial litigation in family
court.}

\noindent{}Grant concluded with plaudits for the author's vision, dedication, and
commonsense approach to the protective mother crisis:

\vspace{-.5em}
\prQuote{Dr. Rosen has performed a huge service by focusing on those of us
who feel numbed by our own inability to protect desperate children and
non-offending parents from the lies of lawyers and psychologists who
have reduced them to a profit center. She [the author] concludes by urging Congress to use its authority and enact CARCO for the District of
Columbia, creating a model for the nation. Federal funding incentives
can be redirected to inspire other states to follow suit and to end the
nightmare that breeds child abuse at family court.\fnm{34}} \pfl{}


\vspace*{1em}
\subsect{15.5}{Protective{\prSd}Mother{\prSd}Issue{\prSd}Gains{\prSd}%
Stardom{\prSd}in{\prSdd}TV{\prSdd}Miniseries\\
Now{\prSe}Streamed{\prSe}on{\prSe}Hulu}

\noindent{}The media coverage of the protective mother’s plight has mushroomed
in recent years, even reaching a Disney Channel platform. In 2022, FX
aired a five-part docuseries, \textit{Children of the Underground}, which would
be subsequently streamed onto Hulu, where it remains to date. The
Emmy Award-winning production company, Story Syndicate, had diligently 
put over three years into the research and production of this
spellbinding television miniseries. Executive Producer Ted Gesing had
contacted one of us in the fall of 2019, after reading the first edition of
this book, saying he was horrified that the family courts had descended
into such a state. The documentary convincingly showed the angst of
mothers and the cruelties of the family court system that oppress 
mothers to the point where many have had no choice but to take their children
and go on the run\mdash{}and suffer
poverty, lack of medical care, isolation, food deprivation, and chronic
stress as a result. \pfl{}

This docuseries on the protective mother issue received good reviews
in a number of different media venues: \textit{New York Post}, 
\textit{The Daily Beast},
A\&E's \textit{True Crime Blog}, and \textit{Indiewire}, 
just to name a few. It was praised
on www.rogerebert.com for \qq{captur[ing] the growing pains of understanding mental health, which includes the people and children who
were mishandled in the [court] process \ellThree{.5pt}{-.5pt} [and is] all the more sobering when one realizes how certain parts of the family court have not
changed, or that the same problems can arise.}\hspace{1pt}\fnm{35} \pfl{}

Yet, some protective mothers and advocates had mixed feelings
about the docuseries. They did not take issue with its content. On the
contrary, they found this docuseries\mdash{}in which mothers and (now)
grown children had so poignantly and persuasively told their stories\mdash{}
very compelling. However, they had thought the gravity of the injustice and eloquence in the way the story had been told would finally
catapult this issue to the front burner. And they were 
frankly disappointed when that did not happen. \pfl{}

However, looking at this situation more objectively, who can say
that the docuseries won’t have that effect; perhaps if not today, then
tomorrow? As with any social movement, the exact point of critical
mass is hard to predict and, similarly, just as hard to identify when
it actually does occur. In fact, it seems that historians are much 
better able to time-stamp critical developments in any political or social
movement than those living through that epoch. And certainly, much
traction has been built up already. Who would have imagined back in
2005 when this book was first published that someday the topic would
occupy a five-part television miniseries airing on a major network and
subsequently streamed onto a Walt Disney platform?\fnm{36} \pfl{}

%\notePage{lii}

\vspace{4em}

\subsect{17}{Women's eNews Recognizes the Protective Mother Issue}

\noindent{}Under the direction of Executive Director and Editor-in-Chief Lori
Sokol, PhD, Women's eNews, the award-winning nonprofit news service that covers issues of concern to women in the United States and
around the world, has welcomed editorials on the protective mother issue. Publishing on this news service has been a real boon to protective
mothers since Women's eNews is a well-respected global forum for
women to share their views on public policy and social reform. \pfl{}

In February 2023, best-selling novelist Talia Carner wrote a 
compelling piece for Women's eNews on why protective mothers, in some
cases, may be perfectly justified in defying family court orders.\fnm{37} Around
the same time, Women's eNews published Amy Neustein's editorial on
the looming civil rights crisis in the family courts. Neustein described
how mothers are routinely deprived of due process and equal protection
of the laws.  One particular injustice described by Neustein is a pattern of
judicial orders that prevent mothers from nursing their infants. She 
offered examples of mothers being court-ordered, under threat of jail and
fines, to express their breast milk into a bottle so that the father could
have generous and unrestricted long weekends and overnights alone
with the infant.\fnm{38} In many of these cases, the fathers
had been absent during the pregnancy and birth. They do not pay child
support either. Yet, they demand the infant spend solitary time with
them, even though it disrupts the infant's nursing schedule and the
mother\hspace{-.5pt}-\hspace{.5pt}child bond. \pfl{}

Women's eNews published a number of subsequent editorials from
Neustein. The pieces looked at family court from various angles. These
included a post-mortem examination of a high profile case wherein a
New York mother, herself an attorney and former federal prosecutor,
took her own life after years of abuse by the Westchester Family Court;
the harsh removal of custody from mothers who had failed to speak
eloquently on the witness stand or had become emotional in court upon
facing an abusive and controlling ex-spouse or partner who was trying
to take their children away from them; the spontaneous formation of
the \qq{Mothers' Vote,} an ad hoc voting bloc of protective mothers across
blue states, red states, and swing states who are united in their family
court plight; and the necessity of going outside the judicial branch to
another branch of government when the courts as an entirety have failed
to protect women and children.\fnm{39} \pfl{}

\enlargethispage{5pt}We noted in our first edition that we had spoken at length to Rita
Henley Jensen, then executive director of Women's eNews, about the
protective mother issue. Ms.\@ Jensen was extremely sympathetic, pointing to the daily calls she was receiving from desperate mothers losing
their children in family court. However, since that was twenty years
ago, the protective mother movement had not yet acquired a unified
voice. Few, if any, articles had been published in Women's eNews on
the protective mother plight. Today, the situation is much different. This
preeminent digital news service has emerged as the premier platform for
sharing in-depth commentaries on the family court system. Certainly,
their wide readership in government and in the private sector will gain
insight into what mothers have been experiencing on a daily basis in the
American family court system. \pfl{}\newpage{}

\vspace{3em}

\subsect{15}{Protective Mother Movement Builds up Traction
on the Political Front}

\noindent{}As an example of the kind of traction built up on the political front,
in 2017, a group of activists from across the country, led by Constance
(\qq{Connie}) Valentine, founder of the California Protective Parents Association,
brought one hundred copies of the first edition of \textit{From Madness to
Mutiny} down to Washington, DC,\fnm{40} and left copies on the desks of the
crucial members of the House of Representatives. Soon thereafter, 
Ms.\@ Valentine and her co-activists were successful in getting the House to
pass a resolution to prioritize child safety in custody and visitation cases.
\qq{H.Con.Res. 72,} as it would officially be called, was sponsored by Maryland Congressman Jamie Raskin, along with several of his colleagues,
who recognized the urgency of putting forth a resolution to affirm the
importance of child safety during family court litigation.\fnm{41} \pfl{}

%\notePage{liv}
Although the Resolution passed in 2017--2018, it left open the possibility of taking the matter, sometime later, to its next level. Written
into the Resolution was a suggestion/provision for a full Congressional
investigation of the family courts that have been proven to show a 
menacing, baleful, and sinister effect on 
the welfare of children since they turn innocent 
children over to abusive parents, while forcing 
the non-offending mother to stand by and watch her children suffer,
and sometimes die, from heinous abuse or neglect. Indeed, when that
next step is reached, the family court backlash against protective mothers will finally be given a front-row seat on Capitol Hill\mdash{}and lawmakers
will listen attentively. \pfl{}


In March 2022, President Biden signed into law the \qq{Keeping Children Safe From Family Violence Act} (commonly referred to as \qq{Kayden's Law,} named after a seven-year-old Bucks County girl [Kayden
Mancuso] murdered by her abusive father during a court-ordered unsupervised visit with the child). Meier, Pollack, and colleagues produced
assiduous research to impress upon Congress the need for a special provision to be included in the reauthorization of VAWA (Violence Against
Women Act) so as to protect vulnerable children from injury and death
at the hands of the abusive father. Seven states, led by Colorado,\fnm{42} have already enacted Kayden's Law, in whole or in part.
They include Colorado, Utah, California, Tennessee, Maryland, Pennsylvania,\fnm{43} and Arizona. Connecticut is now close to enactment as well.
Whilst many more states need to join this effort, the states that have
taken the lead clearly demonstrate the power and force of federal legislation in inspiring states to enact laws to protect children in custody battles
from being subjected to injury and death at the hands of a parent with
a known history of spousal abuse and domestic violence.  The incentives to states are
even more forceful when tied to federal funding. \pfl{}

While the effects of VAWA have evidently trickled down to the state
level, individual states have, similarly, on their own initiated legislation in response to tragic deaths of children in their state. For example,
in 2023, several New York senators co-sponsored \qq{Kyra's Law,} which
was named after a two-year-old Long Island girl forced under court order to have an extended visit with her violent father in Virginia who
%\notePage{lv}
fatally shot her while she slept. \qq{Kyra's Law} emphasizes the importance
of giving sufficient evidentiary weight to claims of domestic violence
when determining the custodial and visitation arrangements for a child.
The bill stresses that keeping a child safe from a parent with a history
of violence should be the first priority of the family courts. Although
many of these state bills, such as Kyra's Law, have not yet been able to
gain the full support of both state-legislative houses needed for passage into law, the passionate commitment of a handful of outstanding
legislators gives reason for hope that such bills will eventually make it to the
floor. \pfl{}

On the whole, this second edition of \textit{From Madness to Mutiny} is being
written at an auspicious juncture in women’s history. We are witnessing
a collective uprising as mothers, advocates, and scholars unanimously
urge an end to violence (both personal and institutional) against women
and children in the family courts across the globe. As demonstrated in
the previously discussed United Nations Human Rights Office of the
High Commissioner study, the family courts have become emblematic
of the ghastly violence against women and children at every phase of
the judicial process. And we see that the mothers are determined to no
longer stand silent in the face of such atrocities\mdash{}both now and in the
future. \pfl{}


\prPart{PART{\prSf}III\mdash{}SEVERE{\prSf}MISTREATMENT%
{\prSf}OF{\prSf}BLACK,{\prSf}INDIGENOUS,\\AND{\prSg}PEOPLE{\prSg}%
OF{\prSg}COLOR{\prSg}MOTHERS{\prSg}IN{\prSg}THE{\prSg}FAMILY{\prSg}COURT}

%\subsect{16}{PROTECTIVE MOTHERS SUFFER TWO FOLD IN THE FAMILY COURTS}

\noindent{}In the new material provided as addenda to a number of the chapters
from the first edition, we examine in detail the mistreatment meted
out to Black, Indigenous, and People of Color (BIPOC) mothers in the
family courts. Although we had presented some compelling material in
the first edition showing, for example, the jailing of a BIPOC mother
from Brooklyn thanks to a capricious ruling of the family court judge,
we hadn’t examined the family courts and the child welfare system
from a racial perspective per se. Now that we have, we can show that
in the closed and claustrophobic world of family courts racism 
flourishes, adding undoubtedly to the severe mistreatment inflicted upon
protective mothers. \pfl{}%\newpage{}


\vspace{3em}

\subsect{15}{Black, Indigenous, and People of Color Mothers Suffer Racial
Discrimination on Several Fronts}

\noindent{}First, BIPOC mothers are slandered by name-calling in the family
courts. In addition to the routine psychiatric name-calling directed
at protective mothers (\qq{disturbed,} suffering from \qq{personality disorders,} \qq{narcissistic,} \qq{grandiose,} etc.) women of color shoulder the extra
burden of having racial epithets hurled at them. For example, in the
Brooklyn Family Court, a Latino mother, who was also a police officer, pleaded with the family court judge for the protection of her
five-year-old girl from alleged acts of sodomy at the hands of the father. 
According to the girl's report, her father had forced himself on her
during an overnight visitation and the child came back hysterical, crying how she \qq{almost choked.}  But when the mother entreated Brooklyn Family Court Judge
Leon S. Deutsch (now deceased) to protect her child, he became enraged. 
calling her (in open court) \qq{an hysterical Hispanic.} \pfl{}

Second, BIPOC mothers are more likely to be incarcerated for contempt of court than their Caucasian counterparts. Interestingly, when comparing the incidence of mere
\qq{threats} of contempt of court (as opposed to acting on those threats
by jailing mothers), we did not find a higher incidence of threats made
against the BIPOC mother. For example, both demographic groups of
%\notePage{lvii}
mothers are equally threatened by family court judges with contempt of
court if they would so much as \qq{speak} to anyone (and sometimes this
includes parents and/or siblings) about their family court proceeding.\fnm{44}
But, when it comes to actually jailing mothers, BIPOC mothers are incarcerated more often than Caucasian mothers, which is very much in
keeping with the statistics across the general population that show a much
higher percentage of incarceration among people of color. \pfl{}

Third, when the family court strips a BIPOC mother of custody and then
orders supervised visitation between the mother and her child (\qq{supervision} is used so as to prevent the mother and child from speaking
about abuse, sexual and otherwise, committed by the father, and it also
functions as a hedge against the child making new 
disclosures of abuse), the mother is often forced to see her
child at the local \textit{police station}, instead of at a court-contracting visitation setting, as described in the next section. Interestingly enough, a
judge's ordering of supervised visits at the police station is not always
contingent on whether or not the BIPOC mother has the 
money to pay the visitation supervision costs when visits are 
held in a private setting. In
studying the family court backlash against protective mothers for 
four decades, we have seldom seen Caucasian mothers ordered into
a police setting to see their children. Such arrangements are appalling,
and they naturally have injurious effects on BIPOC mothers and their
children. \pfl{}

When children see their mothers in the police station, they naturally
develop frightening associations with those visits. As they sit in this 
sterile setting, surrounded by men and women in uniform\mdash{}many of them
carrying guns\mdash{}they often cannot relax on the visit with their mother.
BIPOC families have grown to fear police since they are often the target
of police harassment (sometimes tragically leading to injury and death),
whereas Caucasian families are more likely to view police as people to
whom they can appeal for help. By not giving a BIPOC mother any 
alternative other than a police station for seeing her child is injurious to
the mother and child for many reasons. At the minimum, such visits
cause trauma especially for people of color, who try their very best to
avoid contact with police altogether, and for good reason. Moreover, a
child seeing that their mother must be supervised in a police setting will
understandably think their mother did something terribly wrong, 
otherwise why would she need to be monitored at a police station?
In short, in as much as
it is common for fear of arrest to constantly hover over the
BIPOC population, arranging visits for the mother at a police station is
wholly destructive to the 
mother\hspace{-.5pt}-\hspace{.5pt}child relationship. Any angst that the
mother feels sitting in a police station can be easily transferred to the
child, making the visits uncomfortable (and often traumatic) for both
mother and child. \pfl{}

We followed a case in which a Bronx Family Court judge had ordered
the mother to have her visits with her four-year-old daughter (who 
allegedly was sexually abused by the father) at a police station. The mother
had presented medical evidence to the court to substantiate the abuse,
but the judge refused to listen. The child’s pediatrician had backed up
the credible evidence of abuse. And although the child herself gave 
narrative accounts of how she was abused, neither the law guardian nor
the judge heeded the mother’s pleas to protect her child. The judge, 
instead, stripped the mother of custody and transferred the custody of
the very young girl to the father. The mother held an executive 
secretarial position at a Fortune 500 company. So, she could have paid
for visitation supervision at a private setting, but the judge forced this
BIPOC mother to see her child at the police station instead. This was
the only venue the judge would allow for the mother to have visitation
with her child. The mother remarried, but continued to fight for her
child. \pfl{}

Sadly, her new husband\mdash{}a lawyer\mdash{}was unable, despite 
repeated filings in New York's Appellate Division on her behalf, to alter the dra-
conian visitation arrangement at the police station, let alone to overturn
the custody award to the allegedly abusive father. The mother, worn
down by repeated failures in her efforts to seek justice, was eventually 
unable to continue the fight. This took a tremendous toll on her
marriage, which eventually fell apart. She moved in with a relative in a
southern state, unfortunately far from her daughter. When we asked her
husband about whether racism had played a role in the bizarre visitation
rulings, he didn’t hesitate for a moment. He exclaimed, \qq{For 
sure \ellThree{0pt}{0pt} and throughout the entire case!} \pfl{}

Fourth, whereas family courts are weaponized against protective
mothers, those who are BIPOC are sometimes precluded from the court
process altogether, and that can leave the mother with no redress at all.
For example, a Black special education teacher living in Georgia was
%\notePage{lix}
told last year that she could not even have a hearing before a judge because there are \qq{no available transcribers.} This was untrue\mdash{}the same
Georgia court had made itself available to hear cases brought by Caucasian mothers. Instead, this special ed teacher was shunted to a mediator in lieu of
a judge, as the court fell back on its excuse that there were no court
reporters available. The mediator rushed the case ahead, 
sending her sexually abused seven-year-old daughter 
into the father's custody, and subsequently 
terminated the mother's visitation privileges with her daughter
as well. \pfl{}

The mediator put an end to the mother’s visits with her daughter based
on nothing more than hearsay offered by the father. The mother could
not rebut this evidence because she had been denied access to the courts.
A year later, the mother has still been denied a hearing before a judge,
and she still cannot see her child. The child suffers from a congenital 
disease that has required multiple hospitalizations. The mother, who had
been at her daughter’s side throughout all these hospitalizations, is now
barred from seeing her daughter. Tragically, soon after the termination
of the mother's visits, the child required yet another hospitalization. Yet,
when the child was admitted to the hospital for blood transfusions and
dialysis, the BIPOC mother was forbidden to see her there. Speaking to
one of the authors, she said she feared \qq{she'd
never see her daughter again on earth.} \pfl{}

%\vfill\pagebreak{}
 
\prPart{PART IV\mdash{}UNSAVORY \qq{BUSINESS MODEL} TAKES HOLD IN THE FAMILY COURTS}

\subsect{16}{The New Economic Realities of the Family Courts\mdash{}and How
It Affects Protective Mothers}

\noindent{}The new business paradigm found in the family court system today
forces mothers into accepting services at court-contracting agencies,
such as those that provide visitation supervision. For context, mothers have traditionally been asked to pay for services ordered by the
%\notePage{lx}
court. For example, judges for years have been ordering mothers to pay
for a court-appointed custody evaluation, costing on average \$15,000
for the evaluation and the expert's report made to the court. Today,
however, in addition to the fees for a custody evaluation, 
mothers must pay, for a visitation center facility to monitor their visits with their children. Often paired together with the costs of supervised visitation is
the judge's mandate that, as a precondition for visitation privileges, the
mother must go into \qq{therapy} with a court-approved mental health
practitioner, whose sole purpose is to disabuse the mother of 
her \qq{misperceptions/delusions} about the abuse her children 
have suffered at the hands of the father. \pfl{}

Often these costs are so prohibitively expensive for the mother that
she may stipulate to custody and visitation arrangements\mdash{}undoubtedly
prejudicial to her\mdash{}merely to be relieved of paying such costly 
visitation supervision fees. As a matter of fact, the way that many
of the family court visitation
programs legitimate their charging of protective mothers to see their
children is that they claim to maintain a (discretionary) sliding scale.
However, in practice, while these visitations centers subsidize the costs
for low-income families, such as in cases of inadequate guardianship
where children go hungry, miss school, and receive no medical care,
they charge the middle-class protective mother at full rate. \pfl{}

%\newpage{}

\vspace{2em}

\subsect{15}{Court-Ordered Visitation Used to Prevent the Child's Disclosures
from Investigation}


\noindent{}When the courts order supervised visitation for the mother, the \qq{supervision} is used for two purposes: first, to prevent the mother and
child from speaking about sex abuse committed by the father; and, second, to act as a deterrent against the child initiating discussion with the
mother in making \textit{new} disclosures of abuse. 
Even if a child discloses abuse 
during a supervised visit, the court-appointed visitation monitor is likely
to deny that the disclosure took place, because visitation centers get all
of their business from family courts.
Consequently, they are reluctant to
antagonize a judge who has already labeled the mother a 
\qq{parental alienator} and is hence closed-minded about the abuse. 
As a whole, such a
court \qq{business model} is, for obvious reasons, deeply flawed and counterproductive at best, but inherently deceitful at worst.  What follows is an elucidation
of just how well these visitation centers\mdash{}described by
critics as \qq{a cottage industry}\mdash{}are faring financially in today's
family courts.\fnm{45} \pfl{}

\vspace{2.5em}

\subsect{15}{Visitation Centers Impose Undue Financial Hardship on Mothers}

\noindent{}First, the visitation supervisor often works for just a fraction above
minimum wage (\$16 an hour) to a maximum of \$50 an hour.
Nonetheless, the mother is charged about \$140 to \$160 per hour for the
visit with her child, appropriating the bulk of the hourly fee to 
the visitation center itself. Many supervised 
parents are prohibited from engaging relatives
or friends to supervise their visits, which would certainly relieve most
mothers of an onerous financial burden of paying strangers to supervise 
their visits. The reason given for this is that
the court-contracting
visitation centers must \qq{monitor} and \qq{regulate} these visits, wherein visitation
supervisors are obligated to write regular progress reports to the court.
If one considers, on average, that mothers will be given four hours of
supervised visitation per week, sometimes divided into two visitation
sessions, the weekly costs add up very quickly. \pfl{}

Second, most court-contracting visitation centers require \qq{orientation} fees to be paid by the supervised parent up front before the
visits can be arranged by the visitation facility. Orientation fees can
range anywhere from \$150 to \$450, depending on the region of the
country where the visitation centers are located. During these orientation sessions, the mother is told that her behavior will be closely
%\notePage{lxii}
watched. They are cautioned that notes will be taken on their visits
and that the visitation supervisor will be reporting regularly to the judge. What
is at stake is certainly not an appraisal of how well the mother relates
to the child (no one has ever accused protective mothers of being
deficient in nurturing and loving toward their children), but whether the 
mother \qq{talks against} the father to the child or whether 
the mother gets \qq{angry} at the father if the child should make a new disclosure of sexual
abuse. \pfl{}

\vspace{2.5em}

\subsect{15}{Visitation Centers Impose Punitive Conditions on Mothers
and Children}

\noindent{}Mothers are not only restricted in what they can say to their children,
but they are also forbidden from engaging in normal mother-child 
activities. For example, they are warned not to bring any food or drinks
during their visits with their children, as that can create a burden for the
visitation facility in cleaning up afterwards. \pfl{}

In a recent Manhattan Family Court case, the protective mother was
cautioned during her orientation session with the visitation center not
to bring any gifts for her five-year-old girl because the law guardian
and the father's attorney felt that would be a way of \qq{winning} the
child's loyalty to the mother, which would have been prejudicial to the
father. In fact, in this case, once the visitation supervisor showed sympathy for the mother and urged the court to allow the child to spend
overnights with the mother, unsupervised of course, the visitation supervisor was removed from the case and replaced by another supervisor,
who was tepid toward the mother at best. Similarly, in a Brooklyn case,
the mother was punished with a suspension of her visits because she bent
her head down and whispered \qq{I love you} in the ear of her precious
six-year-old daughter sitting on her lap. When her visits were restored months later, the judge issued an order forbidding the protective
mother from \qq{whispering} to her child; to enforce that order, the visitation supervisors never allowed the young child to sit on the mother's lap
again. \pfl{}
%\notePage{lxiii}

\vspace{1em}

\subsect{15}{\makebox[\textwidth][s]{%
Mothers Must Pay for Court-ordered \qq{Therapy} to Visit Their}\\
Children}
\vspace{-.5em}

\noindent{}Unfortunately, the costs to the protective mother do not end here. As
previously mentioned, in addition to paying for her supervised visitation, and for the orientation at the visitation center, the mother must
pay for \qq{therapy} as well. 
Such court-ordered therapy sessions serve
as a condition that the mother must meet to have visitation with her
child. These are usually weekly sessions and sometimes they are twice a
week. The therapist is then ordered to make regular reports to the court
and to the visitation center on the progress made in those sessions with
the mother\mdash{}effectively using those sessions to disabuse the mother of
the fact that her child had been sexually abused by the father: in other words,
\qq{gaslighting} the mother until she starts to doubt her own sanity. The average cost of those \qq{therapy} sessions paid by the mother is \$150 to \$200.
And those therapy sessions must continue for as long as the mother has
visitation with her child. When the cost of therapy is added to the visitation
supervision costs, the protective mother is burdened with an average of
\$750 to \$800 per week to see her child. \pfl{}

\vspace{3em}

\subsect{15}{Family Court Arrogated by \qq{Ransom} Business Model Spanning
in Many Directions}
\vspace{-.5em}

\noindent{}Taking a child away from a mother without justification and
then forcing her to pay to see her child is tantamount to making a mother pay
ransom for her child. As previously mentioned, this is a relatively new
phenomenon; it had not shown signs of being systematic and routinized
at the time we performed our research for the first edition. Similar to the
business-driven visitation supervision racket, we have detected another
scenario in family court in recent times that likewise follows the ransom
paradigm. To wit, children are wrenched from their mothers without a
hearing and without any evidence before the court. The mothers are then
ordered to pay for an expensive custody evaluation. Having already lost
custody, the mothers are certainly in no position to refuse\mdash{}what other
choice do they have? \pfl{}

%\notePage{lxiv}
This practice of wrenching first and then later performing an
evalua\-tion\mdash{}in essence, a \qq{fishing expedition} to justify an unwarranted
re\-moval\mdash{}bears the indicia of ransom. Let's look at this process more
closely for a moment. \pfl{}

In the past, children primarily stayed with the mother\mdash{}the primary
caregiver\mdash{}until the court-appointed custody evaluator submitted their
report to the court. Only after a full trial wherein the evaluator testified
and was cross-examined did the judge then make a custody 
determination in the form of a written decision, often in the mother's disfavor if
the custody evaluator held tendentious views about women. But today,
the courts instead remove the child from the mother first, and then
order the mother to undergo a court-ordered forensic evaluation. The
fees charged for this court-ordered exam are extravagant because the
so-called neutral experts nesting in the courts depend on these fees for
income. Certainly, had the mother not already lost custody, she might
have been in a position to object to such a costly custody evaluation of
her parenting skills. But when her child has already been taken from
her, the mother, aching for the return of her child, will pay any amount
to get her child back and that includes a court-ordered evaluation at an
extravagant rate. \pfl{}

The court's appointment of a custody expert only \textit{after} the mother
had already lost custody begs the question of how the judge can remove
the child from the mother in the first place as no court-ordered 
evaluation report had yet been before the court? Certainly, these are not
mothers who are active drug users (most don't even have a history
of drug use), or those with DUIs or arrest records of any kind, which
might be justification to remove a child from its mother even before a
full-fledged parental fitness evaluation was performed. Instead, where is
the evidence before the court to warrant taking away the child from the
mother if the custody evaluation had not yet been performed? \pfl{}


\vspace{3em}

\subsect{15}{%
\makebox[\textwidth][s]{Broad Socioeconomic Ramifications of the Ransom Business}\\
Model in the Family Courts}


\noindent{}This \qq{ransom} business model that exists in the court system today, in
which payments are exacted from mothers after their children are
seized without legally required hearings to justify their actions, has
broad economic and social consequences that are important to note.
First, middle-class mothers can ill afford the court ransom fees, which
are heaped on top of the expenses that mothers pay when they lose 
custody to the child's father. They pay, on average, \$2,000 a month for child
support, and to that there are the added costs of healthcare insurance
and tuition for the child. Besides, the costs to protective mothers may
be even greater at times. \pfl{}

For example, in a Long Island, New York, case, the father accused of
sex abuse won custody of his two young children and then asked that
the mother pay for \qq{child care services} on top of all the other expenses
she was forced to pay. How ironic: when the children were with the
mother, her parents babysat while she went to work, and they certainly
didn't charge for their services. But once the mother raised a sex abuse
allegation and lost the custody of her children, her own parents were
then forbidden to see their grandchildren. As a result, they were now
prevented from providing \qq{free} child care for the young children,
which they had been doing all along. Interestingly, although the children
had a set of paternal grandparents living nearby, the father wanted the
mother to pay out-of-pocket for child care when he was unable to care for
them. \pfl{}

\vspace{-1.1em}
Child care expenses, sometimes seen by the courts 
as a legitimate expense to level against the protective mother, drain the mother even
further. This example points to the fact that by eviscerating mothers
from the lives of their children, the entire family ecosystem is disrupted:
maternal grandparents who had hitherto provided free child care so that
their daughters can go to work are now eclipsed from the picture, and
fathers must find hired help to stand in place of the child care provided
by maternal grandparents. Moreover, can the love of grandparents
be replaced with hired help? \pfl{}

Adding to the disruption of the family ecosystem and family stability
is the fact that protective mothers bankrupted by the family courts must
often depend on their parents for financial assistance. In preparation
of this second edition, we interviewed many protective mothers who
have explained how their own parents depleted their retirement funds
to help them shoulder the enormous costs imposed by the ransom 
business model ruling the family court system today. And if these mothers
file for bankruptcy protections, as some of them have, they are still not
protected by bankruptcy laws. This is because the bankruptcy process
protects a mother from her creditors but does not shield her from new
costs. The harsh reality of this distinction is illustrated by the case of a
Houston, Texas mother who filed for bankruptcy but could not persuade
the family court to waive the costs of her supervised visits, essentially
rendering her bankruptcy protection moot. Sadly, she ran out of funds
and lost all contact with her children when she was no longer able to pay
for the supervision of her visits with her children. \pfl{}

So, the mother is left with little choice: either she gives in to the
extortion attempt, paying the ransom to see her children, 
or suffers estrangement from them. Since bankruptcy proceedings do not absolve
the financial burden placed on a mother by the family courts (and this is
for good reason because if one could discharge child support obligations
in bankruptcy proceedings, few fathers would pay support) at the
end of the day the protective mothers, impoverished by family court 
litigation, child support, visitation supervision, and therapy sessions, are
faced with a Hobson's choice. Either they impose a financial burden on
their own parents to exhaust their retirement accounts, placing them in
a dangerously weakened financial position, or they abstain from asking
for help and accept never seeing their children again. \pfl{}

As we can see, the sociological consequences of a criminally dysfunctional family court system creates immense suffering and financial
insecurity for large segments of the population, most notably retirees
who can ill afford to deplete their nest eggs. What this shows is that family court venality has multigenerational repercussions\mdash{}for the mother,
the child, and the grandparents.
The family court has degenerated into a toxic system that
attacks multiple generations\mdash{}weakening the family structure and 
society, and inevitably promoting a shady business model that preys upon
battered women and their abused children. \pfl{}%\newpage{}

%\notePage{lxvii}

\prPart{PART V\mdash{}NARRATIVE CASE HISTORIES AND COURT
RECORDS\raisebox{2pt}{:} A LONGITUDINAL VERSUS A 
CROSS-SECTIONAL STUDY}



\subsect{14}{\makebox[\textwidth][s]{How the Second Edition is Organized, and How It is Distinguished}\\
from the First Edition}

\noindent{}Before discussing the differences between the first edition 
and the second edition, we must point out the following detail.
As with the first edition, we have in the second edition, likewise, prioritized the confidentiality of the families whose cases we analyze and
dissect. Thus, we refrain from citing names of parents and children
when presenting narrative data. In addition, when quoting verbatim
from court documents we have redacted the names of parents and children. No docket numbers are given either, even though in some states
these records are publicly accessible, because our purpose is to protect
the privacy of parents and children where we can. As with the first edition, we have
copiously documented each case so that, in the 
event that either congressional leaders and/or the Department of Justice should decide to
pursue an investigation of the scandalous family court system, we will
be able to direct them to the mothers whose cases are profiled in this
book, as well as to other cases in our database that have not been included in this volume. We find that the court records, when analyzed in
painstaking detail, as we have tried to do in this book, provide a
sufficiently compelling case for a federal investigation. The harrowing
picture presented in mothers' narratives, as relayed to us, is sadly corroborated by the court record itself. This is why we have consistently
pointed to the record as an instantiation of the madness prevailing in
the family courts. Although this process of reviewing extensive court 
files is generally time-consuming, it is important since it strengthens
the terrifying narratives provided by the mothers who are caught in the
vortex of family court litigation. \pfl{}
%\notePage{lxviii}

Our approach in the second edition, however, is distinguished from
that of our first edition in several ways: \pfl{}

First, rather than instantiate, as we did in the first edition, the 
manifold indicia of family court madness and the resultant parental mutiny
by pointing broadly to as many protective mother cases illustrative of
this phenomena as possible, we concentrate instead on fewer exemplary
cases but analyze them in greater detail, and over a longer period of time.
We decided on the longitudinal approach because our first edition, using
a substantial number of cases for illustrative purposes, had already laid
out the themes pervasive in the family courts (themes which still persist
today); therefore, there was no need to repeat that process. Although for
this new edition we hone in on fewer cases, albeit in much more depth,
those we chose to present in this book from the many that we had 
reviewed should not be considered an anomaly. On the contrary, they are
truly representative of the egregious abuses present in the family courts
today. The longitudinal cases presented in this book are, nevertheless,
drawn from various regions of the country to underscore that the 
phenomena they represent are not exclusive to one particular geographic
region. \pfl{}

Overall, our presentation of a few cases in depth, as opposed to
broader analyses of multiple cases, allows us to narrate a given 
case longitudinally and show, for example, how child protective services can
become increasingly more aggressive toward the protective mother as
their case progresses and, in particular, increasingly hostile toward a
protective mother who is a woman of color. Our lengthy case-narrative
approach also allows us to show, for example, how law guardians have
come to conduct themselves in such a reckless manner that they 
arrogantly ignore well-documented criminal histories of fathers over time,
which have a direct bearing on the safety and well-being of a child. \pfl{}

Second, although in preparation of both our first and second editions of this work we immersed ourselves in the study of full court
files\mdash{}transcripts, orders to show cause, forensic reports, motions and
cross-motions, child protective service petitions/case summaries/safety
plans, decisions, appeals, and so forth\mdash{}for our second edition we had
the opportunity to examine the pleadings from the \textit{federal} courts as
well, since in recent years many more mothers have brought civil
rights actions in federal courts in a generally futile attempt to seek remedies
%\notePage{lxix}
they were not able to obtain in the state courts. Even though most, if
not all, mothers had their federal cases dismissed because of collateral estoppel issues and related \qq{abstention} doctrines (e.g., \textit{Younger v. Harris};
\textit{Rooker-Feldman}; \qq{domestic relations exception}) that prevent federal
courts from intervening when a matter is before a state court, their
federal lawsuit papers nonetheless offered \qq{gold mines} of information,
demonstrating the procedural defects/violations pervading their family
court cases. \pfl{}\newpage{}

Third, in researching protective mother cases, both for this new 
edition and for the prior one, we employed a qualitative (rather than a
quantitative) research method. Our goal for this current work was to
use close analytical inspection to identify \qq{patterns of persecution} of
protective mothers that could later be made part of a much larger 
\textit{quantitative} study of such phenomena. That is, by dissecting the family court
institution and its disastrous effects on mothers and children with the
kind of methodological scrutiny we employed, we were able to shed light
on the sometimes obscure structure of a badly malfunctioning system. \pfl{}

Fourth, although we instantiated family court madness in our first
edition by drawing from a broad demographic mix, which naturally 
included cases in the BIPOC population, in our new edition we \textit{specifically}
analyze the effects of racism on the mistreatment of mothers by judges
and by child protective services. \pfl{}

For example, we provide a comparative analysis of how the same child
welfare agency and, moreover, the same medical expert consulting for
that agency, demonstrated a harsher attitude toward the BIPOC mother,
in which every statement she made or action she took were somehow
used against her. She even felt intimidated by her own attorney (who
she eventually fired), fearing his reaction to her taking her four-year-old
daughter to the hospital emergency room after an incident of alleged
sexual assault during an overnight visitation with the father. We formed
our analytical conclusions based on our examination of court records
and child protective service case files and by paying close attention to
the narrations of BIPOC mothers and, sometimes, their family members
as well. All in all, racial discrimination adds another layer of insult
for women of color who are already fighting a blatantly gender-biased
family court and child welfare system. \pfl{}

In looking at the racist elements in the cases we studied, we were surprised to see a dearth of professional literature on how racism plays
an important role in protective mother cases. That is, racial factors in
protective mother cases have not served as an avid topic of study, although it certainly deserves rigorous examination. Furthermore, even
though there are a few women of color who have participated in protective mother support groups, they are nonetheless underrepresented
in such groups. Thus, we use the opportunity in writing this second
edition to give proper attention to the racial factors that are essential in studying family courts' mistreatment of women and children.
Basically, when a system descends to its lowest form of human degradation, as evidenced by unrestrained racial bigotry, this must serve
as a clarion call that something is terribly awry in the system writ
large. \pfl{}

We were surprised to see a dearth of professional literature on the 
important role of racism in protective mother cases. That is, racial factors
in protective mother cases have not served as a focus for study, 
although it certainly deserves rigorous examination. Furthermore, women
of color are noticeably underrepresented in protective mother support
groups. For these and similar reasons, we have taken the opportunity
of this second edition to call attention to the racial factors that are 
essential in studying family courts' mistreatment of women and children.
In fact, one can make the argument that racism is a good \qq{litmus test}
for evaluating the severity of system malfunction in general. Because
when a system descends to its lowest form of human degradation, as
evidenced by unrestrained racial bigotry, this serves as an indication
that something is terribly awry in the system writ large. \pfl{}

Fifth, although in our first edition we closely examined how bogus
mental health theories emerged in family law proceedings, this new 
volume, which employs a longitudinal approach to the study of protective
mother cases, gave us the opportunity to study how a mother's 
day-to-day existence is turned upside down by the quack mental health theories
that have come to haunt her. Mothers, engaged in recording, 
chronicling, and carefully documenting the sexual violations of their children,
have seen how the truth about those crimes can be so readily turned
upside down by the junk science diagnoses attributed to them. This is
perhaps the greatest betrayal of battered women by the 
family courts. \pfl{}\newpage{}

What's even worse is that, as we show in this second edition, the un-
derpinnings of these new renditions of \qq{parental alienation} theories are
even more outlandish than their archetypes. Analyzing the variations of
Parental Alienation Syndrome (PAS) popularly circulating in the family
courts today, we show how junk science theories aimed at discrediting
protective mothers have become more insidious over the years. \pfl{}

All things considered, by copiously documenting the institutional 
violence perpetrated against women in the family courts, we hope this
second edition will play a role in the collective effort of protective
mothers, advocates, legal practitioners, and students of family law in 
inspiring a major overhaul of our nation’s family courts. This will require
a concerted effort by the US Department of Justice, the US Department
of Health and Human Services (HHS), and the US House of 
Representatives and the US Senate. 
Inasmuch as state legislatures and state judicial
conduct commissions have tried for four decades to reform the courts
without success, we sorely need the additional force of the federal 
government. The situation is dire: far too many mothers and children are
stuck in the quagmire of venal family courts, and some are paying with
their lives. \pfl{}

The second edition is organized as follows: We provide a detailed
supplement to the core chapters in the second section of the book:
Robed Rage, Lawless Law Guardians, Anti-Social Services, and Mental
Health Quackery. In addition, we provide supplementary material to the
third section of the book which focuses on reforms. All supplementary
material may be found under the subheading \qq{Addendum.} \pfl{}



